\section{Simulation of the posterior distributions}
\label{sec:4.3}

The posterior distributions for each Ia SNe compilation and model are approximated by histograms created using the Metropolis-Hastings version of Markov chain Monte Carlo. The requirements for a reasonable approximation to the true (unknown) posterior are:

\begin{itemize}
    \item Sufficient burn-in time for the chain to reach an equilibrium.
    \item A suitable maximum step size to explore the parameter space.
    \item Initial conditions (a \lq\lq{}zero-parameter model\rq\rq{}) which do not introduce bias into the walks.
\end{itemize}

Ensuring that all these conditions were met without resorting to highly complex code and parallel computing was a non-trivial task involving much trial and error. This constituted the majority of the project: although simulations were started in May, only 4065 hours of CPU time actually produced data suitable for further processing.\footnote{The output files were timestamped so that it was possible to tell the duration of the walk.} The rest of the time was spent discovering the aforementioned problems in practice, rewriting the code and running new simulations to determine whether the issue had been resolved. 

In this section, I examine the problems faced in more detail and outline the methods used to ameliorate them. The convergence length of the Markov chain cannot be predetermined. Although the \mcmc{} formalism guarantees convergence for most PDFs, it places no bounds on the number of iterations required before convergence is achieved \cite{Gregory2005,Dunkley2005,Christensen2001}. Existing methods are empirical: one runs a series of tests for non-convergence, the negative results of which imply---but do not confirm---that the chain has converged \cite{Christensen2001}. These methods, which include using second-order Markov chains to calculate a bound on the accuracy of the estimations and comparing the values from multiple chains \cite{Dunkley2005}, were considered unnecessary for such a small number of parameters (compared to eleven or more in \cmb{} calculations \cite{LewisBridle2002}). Instead, following \cite{Gregory2005}, the burn-in time preceding convergence was estimated at $1\%$ of the total walk length, provided that the posterior produced was smooth. 

After some experimenting to find a smooth histogram of $\Omega_m$ from the Constitution compilation, the initial walk length estimate of $5 \times 10^4$ steps was increased to $1 \times 10^5$ for the smaller-error observational data (Constitution 0.2, Gold 0.2, Union 0.2) and $2 \times 10^5$ for the projected sets (Snap, Snap 0.2) for one-parameter models. The number of steps required to produce smooth posteriors scales with the dimensionality of the parameter space: (approximately) linearly for \mcmc{} methods \cite{Dunkley2005} compared to exponentially in the crude MC scenario \cite{Christensen2001}. In practice, this linear scaling is not achieved \cite{Christensen2001,Dunkley2005,LewisBridle2002}. Therefore a balance between smoothness and running time was estimated by multiplying the one-parameter walk length tenfold for each additionally varying parameter. The exception to this was multiple-parameter models involving the time-varying dark energy equation of state $w_1$, which had a long burn-in time: these walks had to be redone with twice the existing steps for smoothness. Computational efficiency required that the total number of steps in the chain was both sufficiently large to widely explore the parameter space (\cref{fig:fig4.2b}) and sufficiently small that the program runs in a sensible amount of time given the time span available for the project. 

The second problem in the algorithm is the possibility of sampling bias introduced by the the neighbourhood size. This is governed by a maximum step size: the actual steps will be a fraction of this (depending on the random multiplier chosen). If the maximum step is too small, the walk does not leave a small subset of the parameter volume (\cref{fig:fig4.3b}). Consequently, it will show fine detail (\cref{fig:fig4.3a}), but only in a small section of the characteristic width, leading to an incorrect posterior. Conversely, if the step size is too large, the chain (\cref{fig:fig4.3d}) will oscillate between regions of high likelihood (\cref{fig:fig4.3c}) without exploring the parameter space in-between. This results in a posterior which sparsely outlines peaks in the likelihood, but does not give any accurate shape to the rest of the space. Thus the same initial conditions can produce two very different chains if different maximum step sizes are selected. Although a method called parallel tempering \cref{sec:6.1} exists which efficiently covers the parameter space using information from a variety of step sizes, it was too complex to implement in this project. Consequently the walk had to be repeated for each model over a range of different step sizes incremented by factors of two using a for-loop. It became evident that there was no step size which was suitable for all parameters, e.g. the Hubble parameter $h$ required a small step because of the small characteristic width of its prior, whereas the dark energy equation of state parameters $w_0$ (constant) and $w_1$ (time-varying) produced smooth walks with much larger step sizes. The core step sizes were $0.1 \times 2k ,\, k \in {-2, -1, 0, 1, 2}$ and additional sizes were added after the walks from these were examined: if the $0.025$ step was too large then one of $0.0125$ was run; if the $0.4$ step was too small, steps of $0.8$ and even $1.6$ were run.

The final issue to address is the choice of initial conditions for the walk. If the initial conditions were too far from the likely parameter values then a large number of steps would be wasted as burn-in time: the section of the walk in which the chain travels towards high-likelihood regions \cite{Gregory2005}. Consequently the initial conditions chosen corresponded to the concordance ($\Lambda$CDM) model, an accepted estimate to the correct values of the cosmological parameters \cite{Hartle2003}. The values of the concordance model vary slightly from reference to reference: I adopted the 1-significant figure values used in \cite{Hartle2003}: $\Omega_m = 0.3; w_0 = \Lambda = -1$; with the addition of $w_1 = \gamma = 0$ (by definition the concordance model has neither time-variation in the dark energy equation of state nor interaction). I chose to write the Hubble constant parameter as $h/h_{72}$ so that a value of 1 gave $H_0 = 72\,\mathrm{km}\,\mathrm{s}^{-1}\,\mathrm{Mpc}^{-1}$ , the current best estimate for $H_0$ \cite{Riess2009}.

Once all three criteria were met, the Markov chains produced use approximations to the true posterior PDFs. This enabled the smooth walks with suitable step sizes (\cref{fig:fig4.2}) to be selected for the next stage in the project: finding credible regions.
